\section{Terminal images and exact point-field refinements}\label{sec:point-refinements}
This section records the point-field information behind the first-wild comparison. The terminal-image estimate works for every prime at least five. At five we determine the two point-field filtrations exactly and construct a base with the necessary first cyclotomic break. The latter construction concerns the base field alone; it does not construct two groups with unequal characters.

Throughout the first two subsections, the residue field is algebraically closed, $\mu_p\subset K$, $p\ge5$, $e_K=p(p-1)$, and $L/K$ is the tame extension of degree $N=p^2-1$. Put $s=pN$. A two-slope model has integral parameter $0<B<pe_K$ and first-level values
\[
                     \alpha=s-B,\qquad \beta=B/p.
\]
The high line $H$ is rational over $L$: scaling its degree-$(p-1)$ root factor leaves simple nonzero roots over the algebraically closed residue field, which lift by Hensel's lemma. Its quotient character is trivial, since the determinant is cyclotomic and $\mu_p\subset L$. Thus the residual image over $L$ is trivial or has order $p$, with unique fixed line $H$ in the latter case. Twice dividing a low point by $p$ gives third-level value $B/p^5$. These assertions also follow from the Newton polygons in Sections~\ref{sec:quotients} and~\ref{sec:restricted}.

\subsection{A terminal plane and a wild image bound}
We recover the high line and its parameter from a small full division field.
\begin{lemma}\label{lem:terminal-plane-refinement}
Suppose $B=pb$ with $p\nmid b$ and the full $p^3$-division field $F/L$ has degree $p^4$. In a basis with $H=\F_pe_1$, the terminal space
\[
 V_2=\{C\in M_2(\F_p):1+p^2C\in\operatorname{Gal}(F/L)\}
\]
has dimension two. Every $C\in V_2$ has $C_{21}=0$, and the kernel of $C\mapsto C_{22}$ is exactly $\F_pE_{12}$.
\end{lemma}
\begin{proof}
Normalize $w=p^4v_L$. Every primitive third-level point $P$ with $v=p^2P\notin H$ has $w(T(P))=b$. Its valuation over $L$ has denominator $p^4$, so $L(P)=F$. If a wild automorphism $g$ has lower break $t(g)$, expansion in a uniformizer with fixed Teichm\"uller coefficients gives
\[
 w(gx-x)\ge w(x)+t(g),
\]
with equality when $p\nmid w(x)$, and strict inequality when $p\mid w(x)$. Indeed the first monomial $c\pi^a$ changes to first order by $ac\pi^a(g\pi/\pi-1)$; later exponents have strictly larger valuation. Formal addition gives $w(T(P+R)-T(P))=w(T(R))$, because its difference is $T(R)$ times a unit. Consequently
\begin{equation}\label{eq:terminal-point-break}
                      t(g)=w(T(gP-P))-b.
\end{equation}
A nonzero terminal matrix cannot kill any $v\notin H$, since $P$ generates $F$. For $v=(a,1)$, a nonzero $C_{21}$ would make the second coordinate of $Cv$ vanish for exactly one $a$. At that value $Cv$ is nonzero and high, while at the other values it is low. Equation~\eqref{eq:terminal-point-break} would give two different breaks for the same automorphism. Thus $C_{21}=0$.

If also $C_{22}=0$, nonvanishing of $C_{11}a+C_{12}$ for every $a$ forces $C_{11}=0$. The kernel of $C\mapsto C_{22}$ is therefore contained in $\F_pE_{12}$, and $\dim V_2\le2$. The first congruence space has dimension at most $\dim V_2$, by the $p$th-power map. The residual image has order at most $p$. If $\dim V_2\le1$, the total order would be at most $p^3$, contrary to $[F:L]=p^4$. Hence its dimension is two and the asserted kernel is the full line.
\end{proof}

\begin{proposition}\label{prop:terminal-parameter-recovery}
Suppose two two-slope models have isomorphic $p^3$-torsion and common full division field of degree $p^4$ over $L$. If $B_1=pb_1$ with $p\nmid b_1$, then $H_1=H_2$, $B_1=B_2$, and their special characters agree.
\end{proposition}
\begin{proof}
The denominator of $B_2/p^5$ forces $B_2=pb_2$, with $1\le b_i\le e_K-1$. In the first model, the terminal matrices of the preceding lemma have exact breaks
\[
 t_{\rm low}=(p^4-1)b_1\quad(C_{22}\ne0),\qquad
 t_{\rm high}=p^5N-(p^5+1)b_1\quad(0\ne C\in\F_pE_{12}).
\]
If the high lines differed, choose a residual vector outside both. Its translate by $E_{12}$ is high in the first model and low in the second. The general ramification inequality in the proof of the lemma gives
\[
 t_{\rm high}\le(p^4-1)b_2\le(p^4-1)(e_K-1).
\]
But $p^5N=(p^5+p^4)e_K$, so
\[
 t_{\rm high}=(p^4-1)e_K+(p^5+1)(e_K-b_1)>(p^4-1)e_K.
\]
Thus the lines agree. A matrix with $C_{22}\ne0$ now gives $b_1\le b_2$ from the low displacement inequality; $E_{12}$ gives $b_1\ge b_2$ from the high one. Equality follows. The common high-line character $\zeta^{c_i-B_i}$ then gives $c_1=c_2\pmod N$, hence equality of special characters.
\end{proof}

\begin{corollary}\label{cor:wild-image-at-least-five}
For two potentially supersingular height-two groups with isomorphic $p^3$-torsion, unequal special characters and nontrivial residual wild inertia, the common full division field over $L$ has order between $p^5$ and $p^9$.
\end{corollary}
\begin{proof}
Nontrivial residual wild inertia excludes the one-slope case, whose first-level points are rational over $L$. The common residual transvection identifies both high lines. Write the full degree as $p^r$. If $p\nmid B_i$, the value $B_i/p^5$ forces $r\ge5$. If $v_p(B_i)=1$, it forces $r\ge4$; equality is excluded by Proposition~\ref{prop:terminal-parameter-recovery}. If $r\le4$ remains possible, both parameters must therefore be $B_i=p^2k_i$ with $1\le k_i\le p-2$. Equal high-line characters give $c_2-c_1\equiv B_2-B_1\pmod N$. Equal determinants give $c_2-c_1\equiv0\pmod{p-1}$. Hence $k_2-k_1$ is divisible by $p-1$; its absolute value is smaller than $p-1$, so it is zero. The characters again agree, a contradiction. Finally the residual group has order $p$ and the two higher matrix layers have dimension at most four each, giving the upper bound.
\end{proof}

\subsection{Exact filtrations for the exchanged parameters at five}
We determine the filtrations directly from point coordinates; their exact values contain more information than the conductor bounds used earlier.
\begin{proposition}\label{prop:exact-exchanged-point-profiles}
Consider two potentially supersingular height-two models with tame reducible residual inertia and unequal special characters, at $p=5$, $e_K=20$, with parameters $B_1=25$, $B_2=75$, exchanged high lines $\F_5e_1,\F_5e_2$, and an isomorphism of their $125$-torsion. Assume the residual tame weights are normalized as in the proof of Theorem~\ref{thm:tame-five-recovery}. The full image over $L$ is $1+5M_2(\Z/125)$. For $5P_i=e_i$, the fields $E_i=L(P_i)$ are Galois of degree $25$ and have lower breaks
\[
                  E_1/L:(6,336),\qquad E_2/L:(24,234).
\]
The full second division field has the following breaks:
\[
\begin{array}{c|rrrr}
\text{upper}&6&24&66&72\\
\text{lower}&6&96&1146&1896.
\end{array}
\] Its cyclic degree-five subfields have breaks $6,24,66,72$ with respective multiplicities $1,5,25,125$. The cyclotomic quotient $L(\mu_{25})/L$ has break $72$, and $K(\mu_{25})/K$ has break $3$.
\end{proposition}
\begin{proof}
The first congruence space $W_1$ contains both off-diagonal axes and has both diagonal projections nonzero, by Proposition~\ref{prop:opposite-quotient-image} and its proof. Thus each column evaluation is onto $\F_5^2$. The fields $E_i$ are the corresponding quotient fields of the elementary abelian second-level extension, so they are Galois of degree $25$. The choice of lift $P_i$ does not change its stabilizer.

The first-level high and low values are $(95,5)$ and $(45,15)$. Their second-level values are $(19/5,3/5)$ for $P_1$ and $(1/5,9/5)$ for $P_2$. Multiplication by $[E_i:L]=25$ and formal addition therefore give
\[
\begin{array}{c|c|r|r|r|c}
\text{field}&\text{model}&A&D_0&D_1&H\\\hline
E_1&1&95&125&2375&\text{diagonal}\\
E_1&2&15&375&1125&\text{off-diagonal}\\
E_2&1&5&125&2375&\text{off-diagonal}\\
E_2&2&45&375&1125&\text{diagonal}
\end{array}
\]
Here $A$ is the coordinate valuation, $D_1$ the displacement on the indicated line and $D_0$ the displacement off it. Every displacement is below the absolute index $12000$.

The tame residual eigencharacters are $1,\zeta^6$, up to a common quadratic twist. The diagonal, $21$ and $12$ quotient characters consequently have breaks congruent to $0,6,18$ modulo $24$. Indeed a tame lift on a cyclic degree-five quotient has leading multiplier $\rho$ with $\rho^5=\zeta$, and conjugation on a break-$b$ coefficient is $\rho^b$; multiplication by five fixes these three residues. The two distinct tame weights in each point field force two distinct breaks, because a ramification graded piece has a single tame scalar weight. The last line is the direction with larger cyclic quotient break. These statements can also be read from the quotient index formula of \citet[Proposition~3.7, pp.~14--15, and Proposition~3.13, p.~16]{ColmezLocal}.

Write $d_i,o_i$ for the diagonal and off-diagonal cyclic quotient breaks. In $E_1$, an off-diagonal last line would, by Lemma~\ref{lem:opposite-line-conductor} and the first coordinate row, force $o_1=(125-95)/5=6$. The smaller positive diagonal break would be divisible by $24$, an impossibility. The last line is diagonal, and the second row gives
\[
                         d_1=(375-15)/5=72,\qquad o_1<72.
\]
For $E_2$, a diagonal last line gives $d_2=(125-5)/5=24$. An off-diagonal last line gives $o_2=(375-45)/5=66$ and $d_2<66$. Thus $d_2$ cannot equal $72$.

If $\dim W_1=3$, its diagonal subspace is one dimensional with both projections nonzero. The two diagonal functionals then have the same kernel and define the same field, contradicting $d_1\ne d_2$. Hence $W_1=M_2(\F_5)$. The fifth-power inclusion gives $W_2=M_2(\F_5)$, and counting both layers gives the full congruence image. Its nonzero trace modulo $25$ also shows $\mu_{25}\not\subset L$.

We refine the two filtrations using the fixed-coefficient expansion. If $v_5(A)=1$ and $r$ is the first nonzero prime-to-five exponent of a coordinate, then below the absolute index its only possible leading changes at break $b$ are
\begin{equation}\label{eq:first-wild-leading-changes}
                              A+5b,\qquad r+b.
\end{equation}
To prove this, expand each monomial in $\pi(1+u)$, $v(u)=b$. An exponent divisible by five first contributes modulo five in degree at least five; the initial exponent $A$ has order exactly one and contributes in degree five. Every later exponent has greater value. The earliest exponent prime to five is $r$ and contributes in degree one. Binomial coefficients divisible by five produce terms above the absolute index. Unequal leading values cannot cancel. In particular a displacement greater than $A+5b$, but below the absolute index, forces $r=A+4b$; for any smaller break $t<b$ the displacement is then exactly $A+5t$.

For $E_1$, let $t=o_1$ be the smaller lower break. The first coordinate row and the strict inequality for a coordinate divisible by five give $125>95+t$, so $t<30$. Since $t\equiv6\pmod{24}$, it is $6$. If $u$ is the last lower break, $72=(u+4t)/5$, hence $u=336$.

For $E_2$, a diagonal last line would give larger quotient break $24$, smaller off-diagonal break $18$, and lower breaks $(18,48)$. In its second coordinate row, $1125>45+5\cdot48$ forces $r=45+4\cdot48$ in~\eqref{eq:first-wild-leading-changes}. The displacement at break $18$ must then be $45+5\cdot18=135$, contradicting $375$. Thus its last line is off-diagonal. The larger quotient break is $66$ and the smaller diagonal break $t$ is $24$ or $48$, giving last lower break $u=330-4t$, namely $234$ or $138$. In the first coordinate row, $2375>5+5u$ forces the lower displacement to be $5+5t=125$. Therefore $t=24$ and $u=234$.

The coordinate characters on $M_2(\F_5)$ now have upper breaks
\[
                    A_{21}:6,\quad A_{22}:24,\quad
                    A_{12}:66,\quad A_{11}:72.
\]
A nonzero linear combination has the largest break among its active coordinates, because the other terms vanish on the last nonzero ramification group. Herbrand's inverse function gives successively
\[
 6,\quad6+5(24-6)=96,\quad96+25(66-24)=1146,
 \quad1146+125(72-66)=1896.
\]
Normalizing the last nonzero coefficient of a projective character to one gives $1,5,25,125$ cyclic quotients at these four breaks. The determinant modulo $25$ is $A_{11}+A_{22}$, hence has break $72$. It is the cyclotomic character over $L$. A degree-$24$ tame base change multiplies a cyclic wild break by $24$, as in the uniformizer calculation in the preceding section; the break over $K$ is therefore $3$.
\end{proof}

The proposition is a conditional calculation of the point fields and does not assert that the unequal-character pair exists. Its proof determines the ramification filtration from coordinates and quotient fields, without deriving a table from the impossibility theorem.

\subsection{The exact third-point towers and terminal filtration}
We continue the direct coordinate calculation of Proposition~\ref{prop:exact-exchanged-point-profiles}. This supplies an alternative obstruction in the actual determinant quotient, without using the preceding character-recovery theorem.

Keep the mixed-characteristic base with algebraically closed residue field, $p=5$, $e_K=20$, $\mu_5\subset K$, and tame good extension $[L:K]=24$. Assume the entire comparison hypotheses of Proposition~\ref{prop:exact-exchanged-point-profiles}: two whole potentially supersingular height-two groups have unequal finite special characters, tame reducible residual inertia, isomorphic actual $125$-torsion, parameters $B_1=25,B_2=75$, and exchanged high axes $\F_5e_1,\F_5e_2$, with the residual tame characters normalized there. Put
\[
 F=L(M[125]),\qquad F_2=L(M[25]),\qquad
 W=1+5M_2(\Z/125),\qquad U=1+25M_2(\F_5).
\]
Proposition~\ref{prop:exact-exchanged-point-profiles} gives these full actual images. All lower breaks below use the integer valuation on the top field of the named extension.

\begin{proposition}\label{prop:exact-third-point-profiles}
One can choose a common basis $Q_1,Q_2$ of the actual $125$-torsion with $25Q_i=e_i$ and point fields stable under a tame lift. Set $P_i=5Q_i$, $E_i=L(P_i)$ and $E_i^+=L(Q_i)$. Then $E_i^+/E_i$ is Galois of degree $25$, and
\[
\begin{array}{c|c|r|r}
 &\text{lower breaks}&\text{diagonal quotient}&\text{off-diagonal quotient}\\\hline
 E_1^+/E_1&(606,6936)&1872&606\\
 E_2^+/E_2&(624,6834)&624&1866.
\end{array}
\]
The directions are the actual first-level translations on the two residual axes.
\end{proposition}
\begin{proof}
Choose a lift $\tau$ of the tame generator. Its two residual eigenvalues are distinct. Hensel lifting of the simple roots of its characteristic polynomial and the polynomial idempotents
$(\tau-\lambda_j)/(\lambda_i-\lambda_j)$ split the free $\Z/125$-module into two free rank-one eigensummands. Choose $Q_i$ on them. Then $\tau Q_i$ is a unit multiple of $Q_i$, so the point fields are $\tau$-stable. A finite-order choice of $\tau$ is unnecessary.

Let $S_i$ and $T_i$ be the stabilizers of $Q_i$ and $P_i$ in $W$. The full third-level column orbit has $625$ elements, and the second-level orbit has $25$ elements. For $g\in T_i$, write $gQ_i-Q_i=25b_i$ in the common torsion coordinates. The map $g\mapsto b_i\in\F_5^2$ is a homomorphism, because $g\equiv1\pmod5$. The terminal matrices already make it surjective, and its kernel is $S_i$. Hence $S_i$ is normal in $T_i$, proving the Galois assertion and the precise translation marking. We do not assert that $E_i^+/L$ is Galois.

The second-level point values are $19/5,3/5,1/5,9/5$ in the four field/model rows. All are below their corresponding first-level low slopes. Dividing once more therefore divides each valuation by $25$, while the field degree grows from $25$ to $625$. The normalized coordinate values remain $95,15,5,45$, and formal addition gives
\[
\begin{array}{c|c|r|r|r|c}
\text{field}&\text{model}&A&D_0&D_1&H\\\hline
E_1^+&1&95&3125&59375&\text{diagonal}\\
E_1^+&2&15&9375&28125&\text{off-diagonal}\\
E_2^+&1&5&3125&59375&\text{off-diagonal}\\
E_2^+&2&45&9375&28125&\text{diagonal}.
\end{array}
\]
Every $A$ has five-adic order one, and every displacement is below $v_{E_i^+}(5)=300000$. The tame weights are again zero and six for the first tower, and eighteen and zero for the second. Indeed conjugation on a translation divides by the eigenvalue of the source $Q_i$. A cyclic degree-five quotient over $E_i$ has total ramification degree $125$ over $L$, so its leading tame multiplier satisfies $\rho^{125}=\zeta$. Thus the diagonal quotient break is a positive multiple of $24$, and the off-diagonal quotient break is six or eighteen modulo $24$. There are two distinct lower breaks $t<u$, with an eigenaxis as last line.

For a row with coordinate valuation $A$, let $r$ be the first nonzero exponent prime to five in its fixed-coefficient expansion. The first possible displacement at break $b$, below the absolute index, is $\min(A+5b,r+b)$. In particular a displacement larger than $A+5b$ forces $r=A+4b$; at every smaller break $t$ the displacement is then exactly $A+5t$.

In the first tower, an off-diagonal last line would force the larger cyclic quotient break $(3125-95)/5=606$ by Lemma~\ref{lem:opposite-line-conductor}. Thus $t<606$ and $u=3030-4t$. The second model has $28125>15+5u$, so its lower displacement would be $15+5t=9375$, forcing $t=1872$, a contradiction. The last line is diagonal. The second row then gives larger quotient break $1872$, so $u=9360-4t$. The first row has $59375>95+5u$ and forces $3125=95+5t$. Hence $t=606$ and $u=6936$.

For the second tower, a diagonal last line would give larger quotient break $(3125-5)/5=624$, so $t<624$ and $u=3120-4t$. The fourth row has $28125>45+5u$, forcing $9375=45+5t$ and $t=1866$, a contradiction. Its last line is off-diagonal. The fourth row gives larger quotient break $1866$, hence $u=9330-4t$. The third row has $59375>5+5u$ and forces $3125=5+5t$. Thus $t=624$ and $u=6834$. The quoted cyclic breaks follow from $(u+4t)/5$.
\end{proof}

\begin{proposition}\label{prop:exact-terminal-filtration}
The terminal extension $F/F_2$ would have lower breaks
\[
                  7896,\quad11646,\quad769146,\quad937896,
\]
with coordinate directions removed in the order $B_{22},B_{21},B_{12},B_{11}$ for $1+25B\in U$. Outside $U$, the lower breaks in $F/L$ equal the first-layer breaks $6,96,1146,1896$, in the order $A_{21},A_{22},A_{12},A_{11}$.
\end{proposition}
\begin{proof}
Put $K_i=F_2E_i^+$. The terminal group maps onto $T_i/S_i$, so $T_i=S_iU$. Galois correspondence gives $E_i^+\cap F_2=E_i$ and
\[
 \operatorname{Gal}(K_i/E_i)
 =\operatorname{Gal}(E_i^+/E_i)\times\operatorname{Gal}(F_2/E_i).
\]
Both factors are actual Galois extensions of degree $25$. The group $\operatorname{Gal}(K_i/F_2)$ is the $i$th terminal column $U/(U\cap S_i)$. The two column kernels have trivial intersection, so $K_1K_2=F$.

Lower numbering is inherited by subgroups. In the first case the subgroup for $F_2/E_1$ has coordinates $A_{22},A_{12}$, and in the second it has $A_{21},A_{11}$. Their lower and upper break lists are
\[
\begin{array}{c|c|c}
 &\text{lower}&\text{upper}\\\hline
 F_2/E_1&(96,1146)&(96,306)\\
 F_2/E_2&(6,1896)&(6,384).
\end{array}
\]
The new tower upper breaks are $(606,1872)$ and $(624,1866)$. All four breaks in each elementary abelian compositum are distinct. The unequal-break sum rule and upper-quotient compatibility therefore give their union, and inverse Herbrand integration gives
\[
\begin{array}{c|c|c}
 &\text{upper for }K_i/E_i&\text{lower for }K_i/E_i\\\hline
 i=1&(96,306,606,1872)&(96,1146,8646,166896)\\
 i=2&(6,384,624,1866)&(6,1896,7896,163146).
\end{array}
\]
Restricting to the terminal column subgroups retains the last two lower values. Converting each column back to upper numbering gives
\[
 K_1/F_2:\ (8646,40296)\quad(B_{21},B_{11}),\qquad
 K_2/F_2:\ (7896,38946)\quad(B_{22},B_{12}).
\]
These markings are restrictions of the actual maps $g\mapsto gQ_i-Q_i$, so they cannot be changed by an abstract identification of Galois groups.

The two columns generate $F/F_2$. Their coordinate upper breaks are
\[
               7896,\quad8646,\quad38946,\quad40296.
\]
Integrating the indices $1,5,25,125$ yields
\[
\begin{aligned}
 b_1&=7896,& b_2&=7896+5\cdot750=11646,\\
 b_3&=11646+25\cdot30300=769146,&
 b_4&=769146+125\cdot1350=937896.
\end{aligned}
\]

Every nonidentity element of $U$ has lower break at least $7896$. In any nonzero coset of $U$, the quotient average shows that some lift has smaller break, because the quotient $F_2/L$ has last lower break $1896$. Multiplying that lift by an element of $U$ preserves the smaller break, so every lift has the same break. The quotient average then identifies it with the first-layer break. This proves the final assertion. The only ramification inputs are subgroup restriction and the norm-polynomial/Herbrand quotient formulas of \citet[Proposition~3.7, Theorem~3.10 and Proposition~3.13(ii), pp.~14--16]{ColmezLocal}.
\end{proof}

\begin{corollary}\label{cor:determinant-fibre-obstruction}
The preceding two exact filtrations cannot arise from the stated actual torsion comparison. They give a second proof of tame reducible character recovery at $p=5$, $e_K=20$.
\end{corollary}
\begin{proof}
The actual determinant on $W$ has image $1+5\Z/125$, of order $25$, and is cyclotomic over $L$. Fix a nonzero $s\in\F_5$ and its order-five element $1+25s$. Write $g=1+5A+25B$ with entries of $A,B$ chosen in $\{0,1,2,3,4\}$. The determinant condition requires $\operatorname{tr}A=0\pmod5$ and
\[
 \operatorname{tr}B=s-\det A-\operatorname{tr}A/5\pmod5.
\]
Thus every trace-zero $A$ has $125$ lifts in the fibre, with the integer trace carry included.

Among the $124$ nonzero trace-zero first-layer matrices, $100$ have $A_{21}\ne0$, $20$ have $A_{21}=0$ with nonzero diagonal, and four have only $A_{12}\ne0$. Their breaks are $6,96,1146$. For $A=0$, the condition is $\operatorname{tr}B=s$; its terminal coordinate counts are $100,20,4,1$ in the removal order. The complete fibre distribution is
\[
\begin{array}{c|rrrrrrr}
 b(g)&6&96&1146&7896&11646&769146&937896\\\hline
 \#g&12500&2500&500&100&20&4&1.
\end{array}
\]
It has $15625$ elements and weighted break sum $5925000$. The exact quotient-index formula of \citet[Proposition~3.7, pp.~14--15]{ColmezLocal}, after subtracting one, would give the nonidentity determinant quotient element lower break
\[
                  \frac{5925000}{15625}=\frac{1896}{5}.
\]
This is impossible for a normalized discrete valuation: a lower break is an integer. This contradiction uses the exact actual-field tables and determinant fibre, without appealing to a later recovery theorem or to Hasse--Arf.
\end{proof}


\subsection{A base with cyclotomic break three}
We show that the first cyclotomic condition in the preceding proposition can occur for a complete base field.
\begin{proposition}\label{prop:break-three-base}
There is a complete discretely valued field $K$ with algebraically closed residue field of characteristic five, absolute index $20$ and $\mu_5\subset K$, such that $K(\mu_{25})/K$ is cyclic of degree five with lower break three, and $K(\mu_{125})/K$ is cyclic of degree $25$. After a tame extension $L/K$ of degree $24$, the extension $L(\mu_{25})/L$ has lower break $72$.
\end{proposition}
\begin{proof}
Let $k$ be algebraically closed of characteristic five, put $K_0=W(k)[1/5]$, $K_1=K_0(\zeta_5)$ and $\pi_1=\zeta_5-1$. The shifted cyclotomic polynomial is Eisenstein, so $e_{K_1}=4$ and $\pi_1$ is a uniformizer. Set
\[
                  f(T)=(1+T)^5-\zeta_5+\pi_1^4T^2.
\]
Its constant term is $-\pi_1$, all its other nonleading coefficients have positive valuation, and it is monic. Thus it is Eisenstein. For a root $t$, the field $K=K_1(t)$ has index $20$ and uniformizer $t$.

Put $z=\zeta_{25}-(1+t)$. The equation $f(t)=0$ gives
\[
 z^5+5(1+t)z^4+10(1+t)^2z^3+10(1+t)^3z^2
           +5(1+t)^4z-\pi_1^4t^2=0.
\]
The constant coefficient has $K$-valuation $22$, while all four intermediate coefficients have valuation $20$. They lie strictly above the Newton segment from $(0,22)$ to $(5,0)$. Hence every root has valuation $22/5$, and $[K(\zeta_{25}):K]=5$. This extension is cyclic because $\mu_5\subset K$.

Normalize its valuation by $w=5v_K$. For a nonidentity automorphism,
\[
 w(z)=22,\qquad
 \sigma(z)-z=(\zeta_5^a-1)\zeta_{25},\qquad
 w(\sigma(z)-z)=25.
\]
Since $22$ is prime to five, the uniformizer-expansion equality gives lower break $25-22=3$.

The extension $K_1(\mu_{125})/K_1$ is cyclic of degree $25$ and has unique degree-five subfield $K_1(\mu_{25})$. A nontrivial intersection with the degree-five field $K$ would force equality with that subfield, contrary to $[K(\mu_{25}):K]=5$. The intersection is therefore $K_1$, proving the degree and cyclicity assertion. The same tame uniformizer calculation as above multiplies the degree-five break by $24$, giving $72$ over $L$. No formal-group existence is inferred from this field construction.
\end{proof}
